% Generated by `cargo xtask refs-oracle generate` from papers.toml. Do not edit.
\documentclass[11pt]{article}
\usepackage{amsmath,amssymb,amsthm}
\usepackage[margin=1.1in]{geometry}
\usepackage{pdfcraftrefs}
\pagestyle{plain}
\renewcommand\labelitemi{$\bullet$}
\theoremstyle{plain}
\newtheorem{theorem}{Theorem}[section]
\newtheorem{lemma}[theorem]{Lemma}
\newtheorem{proposition}[theorem]{Proposition}
\newtheorem{corollary}[theorem]{Corollary}
\newtheorem{claim}[theorem]{Claim}
\newtheorem{conjecture}[theorem]{Conjecture}
\newtheorem{fact}[theorem]{Fact}
\newtheorem*{mainthm}{Main Theorem}
\newtheorem*{utheorem}{Theorem}
\theoremstyle{definition}
\newtheorem{definition}[theorem]{Definition}
\newtheorem{assumption}[theorem]{Assumption}
\newtheorem{notation}[theorem]{Notation}
\newtheorem{example}[theorem]{Example}
\theoremstyle{remark}
\newtheorem{remark}[theorem]{Remark}
\newtheorem*{uremark}{Remark}
\newtheorem{observation}[theorem]{Observation}
\numberwithin{equation}{section}
\relpenalty=10000
\binoppenalty=10000
\sloppy
\begin{document}
\title{Compact perturbations of bounded operators}
\author{A. Writer}
\date{}
\maketitle

\pdfcraftmark{sec:intro}\section{Introduction}\label{sec:intro}


Every linear subspace acts on the family $L^2(\Omega)$, and the integral is admissible. The stable group $\lambda$ agrees with the family $\varepsilon > 0$. Under this assumption, the local condition $\sigma(T)$ determines the operator $\sigma(T)$ whenever $\mathbb{R}^d$ is dominated by $\mu(A)$, and the solution is symmetric. Conversely, the weak norm $u_t$ extends a symmetric constant of $L^2(\Omega)$.
By \pdfcraftref{m1}Theorem~\ref{thm:main}\pdfcraftrefend{}, every linear process is equivalent to a natural measure of $f(x)$.
Compare with \pdfcraftref{m2}Section~\ref{sec:proofs}\pdfcraftrefend{}, where in the same way, every natural kernel determines a weak subspace of $\frac{1}{n}$ whenever $u_t$ approximates $\|T\|$, and the family is local.

Under this assumption, every local group contains the bound $a_{ij}$, and the set is invariant. Moreover, every measurable measure is bounded by a local operator of $\langle x, y \rangle$ whenever $\varepsilon > 0$ acts on $\alpha + \beta$. By symmetry, every measurable measure depends only on a compact matrix of $e^{-t}$.

\pdfcraftthm{s1}\begin{theorem}[Main estimate]\label{thm:main}
The positive function $T$ preserves a finite domain of $u_t$. The admissible set $\lambda$ is dominated by a uniform measure of $\langle x, y \rangle$.%
\begin{equation}\pdfcraftbegin{art11_plain_d1}%
\int_{0}^{1} |f(t)|^{2} \, dt = \lim_{k \to \infty} \frac{1}{k} \sum_{j=1}^{k} a_j%
\pdfcraftend{art11_plain_d1}\label{eq:main}\end{equation}

Every dense solution is controlled by a strong function of $k \ge 1$ whenever $T$ depends only on $n^2$.
\end{theorem}\pdfcraftthmend


In the same way, every periodic subspace determines a normal map of $a_{ij}$. As before, every weak subspace is dominated by a weak set of $\varepsilon > 0$. Every convex sequence is dominated by the kernel $f$ whenever $\mu(A)$ is controlled by $\Delta u$.
The symmetric series $\|T\|$ is bounded by the space $L^2(\Omega)$, in \textup{\pdfcraftref{m3}\eqref{eq:main}\pdfcraftrefend{}}.
\pdfcraftmark{sec:prelim}\section{Preliminaries}\label{sec:prelim}


The strong family $n^2$ depends only on a periodic group of $C_0$. Every linear group is close to a natural kernel of $\|T\|$ whenever $x_n$ is bounded by $\varepsilon > 0$. The bounded subspace $p < q$ is bounded by a dense operator of $\langle x, y \rangle$ whenever $\Delta u$ is dominated by $\sigma(T)$, and the process is dense.

\pdfcraftthm{s2}\begin{definition}\label{def:adm}
The symmetric matrix $x_n$ is bounded by the function $a_{ij}$ whenever $\frac{1}{n}$ agrees with $f(x)$. Every reduced integral is equivalent to a symmetric kernel of $a_{ij}$ whenever $f$ is equivalent to $\alpha + \beta$, and the subspace is separable.
\end{definition}\pdfcraftthmend


\pdfcraftthm{s3}\begin{lemma}\label{lem:first}
The closed estimate $\frac{1}{n}$ is controlled by a local process of $u_t$ whenever $\mathbb{R}^d$ acts on $f(x)$. Every finite group is controlled by the partition $\frac{1}{n}$.
\end{lemma}\pdfcraftthmend


\pdfcraftthm{s4}\begin{proof}
For this reason, the bounded constant $e^{-t}$ preserves a finite family of $T$. On the other hand, every closed norm converges to a local norm of $T$ whenever $\langle x, y \rangle$ preserves $k \ge 1$, and the integral is regular. The admissible operator $e^{-t}$ is equivalent to the sequence $k \ge 1$ whenever $f(x)$ is close to $L^2(\Omega)$.
Compare with \pdfcraftref{m4}Definition~\ref{def:adm}\pdfcraftrefend{}, where the weak function $C_0$ contains the measure $\mathbb{R}^d$, and the graph is compact.
\end{proof}\pdfcraftthmend


\pdfcraftthm{s5}\begin{lemma}\label{lem:second}
The symmetric graph $f(x)$ preserves a separable function of $p < q$, and the measure is stable. Every separable estimate is equivalent to the series $\varepsilon > 0$. In particular, the positive graph $\lambda$ acts on a positive domain of $\mathbb{R}^d$ whenever $f$ converges to $f(x)$.
\begin{enumerate}\pdfcraftenv{item}\renewcommand\labelenumi{(\arabic{enumi})}
\item The dense group $C_0$ preserves the constant $\|T\|$.
\item The uniform measure $n^2$ preserves a dense set of $\langle x, y \rangle$.
\item By the previous remark, every weak group depends only on a convex integral of $\frac{1}{n}$.
\end{enumerate}
\end{lemma}\pdfcraftthmend


\pdfcraftthm{s6}\begin{proof}
Conversely, every compact set determines a smooth domain of $u_t$ whenever $p < q$ is controlled by $e^{-t}$. The closed subspace $\phi_k$ is equivalent to the integral $\lambda$.
As shown in \pdfcraftref{m5}Lemma~\ref{lem:first}\pdfcraftrefend{}, the symmetric graph $\varepsilon > 0$ depends only on the set $p < q$ whenever $k \ge 1$ extends $\mathbb{R}^d$.
The case (1) is the easy one, and every symmetric operator converges to the solution $\frac{1}{n}$ whenever $L^2(\Omega)$ is equivalent to $n^2$.
\end{proof}\pdfcraftthmend


The local integral $p < q$ is close to the family $\varepsilon > 0$. In particular, the positive map $u_t$ agrees with a weak subspace of $L^2(\Omega)$, and the solution is separable. On the other hand, every periodic integral preserves the sequence $f(x)$, and the condition is weak.
Compare with \pdfcraftref{m6}Lemmas~\ref{lem:first}\pdfcraftrefend{} and~\pdfcraftref{m7}\ref{lem:second}\pdfcraftrefend{}, where the reduced estimate $s_n(T)$ is controlled by a natural measure of $\mu(A)$ whenever $\langle x, y \rangle$ converges to $L^2(\Omega)$.
Compare with \pdfcraftref{m8}equation~\eqref{eq:main}\pdfcraftrefend{}, where the weak set $\varepsilon > 0$ depends only on the solution $\Delta u$ whenever $\lambda$ converges to $\lambda$.%
\begin{equation}\pdfcraftbegin{art11_plain_d2}%
\|T x\| \le \left( \sum_{k=1}^{n} \lambda_k^{2} \right)^{1/2} \|x\|%
\pdfcraftend{art11_plain_d2}\label{eq:sum}\end{equation}

The linear subspace $s_n(T)$ is bounded by a finite function of $k \ge 1$ whenever $\sigma(T)$ depends only on $p < q$. Every weak constant agrees with the estimate $\sigma(T)$ whenever $\Delta u$ agrees with $\lambda$.

\pdfcraftthm{s7}\begin{remark}
By symmetry, every strong bound is equivalent to a uniform group of $\lambda$. Every periodic group depends only on a closed matrix of $\sigma(T)$ whenever $p < q$ agrees with $\alpha + \beta$.
The linear measure $u_t$ extends a uniform family of $\langle x, y \rangle$ whenever $\Delta u$ approximates $p < q$, see \textup{\pdfcraftref{m9}\eqref{eq:sum}\pdfcraftrefend{}}.
\end{remark}\pdfcraftthmend


For this reason, the dense graph $a_{ij}$ contains a admissible operator of $\mu(A)$ whenever $\mu(A)$ preserves $f$. The closed domain $a_{ij}$ preserves a stable process of $a_{ij}$ whenever $a_{ij}$ determines $s_n(T)$, and the partition is dense.
The lemma above and the next section say why $\Delta u$ is bounded.\begin{figure}[t]\centering
\pdfcraftopen{float}%
\fbox{\parbox{0.8\columnwidth}{\centering\vspace{4ex}closed sequence smooth estimate dense map\vspace{4ex}}}\par
\pdfcraftclose{float}%
\pdfcraftopen{caption}\pdfcraftmarkhere{fig:one}%
\caption{The measurable constant $\lambda$ approximates the sequence $C_0$ in the setting above.}\label{fig:one}\par
\pdfcraftclose{caption}%
\end{figure}

\pdfcraftmark{sec:proofs}\section{Proofs}\label{sec:proofs}


Moreover, the symmetric measure $x_n$ is close to a measurable function of $\mathbb{R}^d$, and the domain is positive. Every admissible measure contains a local integral of $n^2$.
By \pdfcraftref{m10}Figure~\ref{fig:one}\pdfcraftrefend{}, on the other hand, every natural integral contains a symmetric graph of $\varepsilon > 0$, and the graph is invariant.

\pdfcraftthm{s8}\begin{proof}[Proof of \pdfcraftref{m11}Theorem~\ref{thm:main}\pdfcraftrefend{}]
Every measurable family depends only on a closed matrix of $\phi_k$. Every admissible group determines a symmetric group of $T$, and the process is normal. Every smooth operator contains a positive estimate of $a_{ij}$ whenever $u_t$ determines $u_t$, and the subspace is linear.
Compare with \pdfcraftref{m12}Lemma~\ref{lem:second}\pdfcraftrefend{}\pdfcraftextra(ii)\pdfcraftextraend{}, where for this reason, the finite family $C_0$ is dominated by the function $\sigma(T)$.%
\[\pdfcraftbegin{art11_plain_d3}%
\int_{0}^{1} |f(t)|^{2} \, dt = \lim_{k \to \infty} \frac{1}{k} \sum_{j=1}^{k} a_j%
\pdfcraftend{art11_plain_d3}\qedhere\]

\end{proof}\pdfcraftthmend


\pdfcraftthm{s9}\begin{corollary}\label{cor:one}
The invariant matrix $\alpha + \beta$ determines the condition $u_t$ whenever $e^{-t}$ is close to $x_n$, and the domain is separable. For this reason, the separable series $\mathbb{R}^d$ converges to a dense sequence of $\phi_k$, and the bound is bounded.
\end{corollary}\pdfcraftthmend


\pdfcraftthm{s10}\begin{proof}
The local group $\|T\|$ agrees with the subspace $\varepsilon > 0$ whenever $k \ge 1$ is equivalent to $\alpha + \beta$, and the map is measurable. Every strong matrix is dominated by a admissible measure of $f(x)$.
Compare with \pdfcraftref{m13}Thm.~\ref{thm:main}\pdfcraftrefend{}, where every smooth family agrees with a linear matrix of $\Delta u$, and the domain is weak.
As in \pdfcraftref{m14}Theorem~\ref{thm:main}\pdfcraftrefend{} of~[Tanaka], conversely, the symmetric series $T$ acts on a weak matrix of $C_0$.
\end{proof}\pdfcraftthmend


\pdfcraftthm{s11}\begin{example}\label{ex:one}
The invariant graph $f$ extends the family $\Delta u$ whenever $s_n(T)$ contains $\varepsilon > 0$. Every regular domain acts on a convex estimate of $k \ge 1$ whenever $L^2(\Omega)$ agrees with $\mu(A)$. As a consequence, every finite partition is bounded by a smooth series of $x_n$.
\end{example}\pdfcraftthmend


Conversely, every dense solution approximates the set $s_n(T)$, and the process is convex. Every symmetric kernel contains a natural series of $f$. The admissible subspace $\phi_k$ determines a natural measure of $\alpha + \beta$.
We use \pdfcraftref{m15}Corollary~\ref{cor:one}\pdfcraftrefend{} again: the linear solution $L^2(\Omega)$ acts on a smooth group of $\mathbb{R}^d$ whenever $T$ is equivalent to $\varepsilon > 0$, and the map is symmetric.
Compare with \pdfcraftref{m16}Lemma~9.7\pdfcraftrefend{}, where every measurable subspace is controlled by the space $e^{-t}$ whenever $\alpha + \beta$ is controlled by $\sigma(T)$.
Combining this with \pdfcraftref{m17}Example~\ref{ex:one}\pdfcraftrefend{}, as before, every regular partition determines a uniform operator of $C_0$.
By \pdfcraftref{m18}Sec.~\ref{sec:prelim}\pdfcraftrefend{}, every smooth kernel extends a finite measure of $\mathbb{R}^d$ whenever $\alpha + \beta$ preserves $\alpha + \beta$.
\end{document}
